%!TEX root = ../../main.tex

\chapter{Stand der Forschung}
\label{ch:stand-der-forschung}

Aufbauend auf den theoretischen Grundlagen ordnet dieses Kapitel die vorliegende Arbeit in den aktuellen Forschungsstand ein. Vorangestellt ist das Vorgehen der Literaturrecherche, das die Auswahl der herangezogenen Arbeiten offenlegt (Abschnitt~\ref{sec:recherche}). Anschließend wird der breitere Einsatz von Künstlicher Intelligenz in der Datenverarbeitung betrachtet (Abschnitt~\ref{sec:ki-datenverarbeitung}). Anschließend wird die Forschung zur Verarbeitung strukturierter und tabellarischer Daten durch \acp{LLM} sowie zur Erzeugung synthetischer Datensätze aufgearbeitet, die die Wahl der Datengrundlage dieser Arbeit begründet (Abschnitt~\ref{sec:llm-tabellen}), bevor verwandte Arbeiten zur Automatisierung von \ac{ETL}-Prozessen diskutiert werden (Abschnitt~\ref{sec:verwandte-arbeiten}). Auf dieser Grundlage wird abschließend die Forschungslücke präzisiert, an der die vorliegende Arbeit ansetzt (Abschnitt~\ref{sec:forschungsluecke}).

\section{Vorgehen der Literaturrecherche}
\label{sec:recherche}

Um die Auswahl der herangezogenen Arbeiten nachvollziehbar zu machen, wird das Vorgehen der Recherche vorangestellt. Gesucht wurde entlang der drei Themenstränge, die sich aus der Zielsetzung ergeben: \ac{ETL}-Prozesse und Datenqualität, der Einsatz von \acp{LLM} auf strukturierten Daten sowie Prompt Engineering. Verwendet wurden dafür englischsprachige Suchbegriffe und deren Kombinationen, unter anderem \enquote{ETL automation}, \enquote{LLM data cleaning}, \enquote{data wrangling foundation models}, \enquote{entity matching}, \enquote{data quality dimensions} und \enquote{prompt engineering survey}.
Als Einstiegspunkte dienten Google Scholar und arXiv, ergänzt um die Vorwärts- und Rückwärtssuche ausgehend von den zentralen Übersichtsarbeiten.

Aufgenommen wurden Arbeiten, die einen der drei Themenstränge unmittelbar behandeln und entweder eine empirische Untersuchung, eine systematische Übersicht oder eine begründete Architekturbeschreibung liefern. Der Schwerpunkt liegt bewusst auf aktuellen Veröffentlichungen, da sich das Feld der \acp{LLM} in Monatsabständen verändert; ältere Arbeiten wurden nur dort einbezogen, wo sie einen etablierten Begriff oder ein etabliertes Verfahren begründen, etwa zur Duplikaterkennung. Der so entstandene Bestand umfasst 46~Quellen aus den Jahren 1996 bis 2026 mit deutlichem Schwerpunkt auf den Jahren 2023 bis 2026, auf die 30~Quellen entfallen, davon 28~Zeitschriftenaufsätze, neun Konferenzbeiträge und neun Online-Quellen.

Für die Einordnung der folgenden Abschnitte ist eine Beobachtung zur Beschaffenheit dieses Bestands wesentlich. Die empirischen Arbeiten zur Verarbeitung strukturierter Daten durch \acp{LLM} erscheinen überwiegend in begutachteten Fachzeitschriften und auf etablierten Konferenzen. Die konzeptionelle Literatur zur \ac{KI}-gestützten Automatisierung von \ac{ETL}-Prozessen, die den in Abschnitt~\ref{sec:verwandte-arbeiten} referierten Erwartungshorizont prägt, erscheint dagegen zu einem erheblichen Teil in Publikationsorganen ohne erkennbar strenge Begutachtung. Ihre Aussagen sind entsprechend als Beschreibung eines Erwartungshorizonts und nicht als gesicherter Befund zu lesen. Gerade dieser Umstand stützt die Anlage der vorliegenden Arbeit, denn er verweist auf einen Bedarf an kontrollierter empirischer Prüfung dessen, was dort konzeptionell in Aussicht gestellt wird.

\section{KI-Einsatz in der Datenverarbeitung}
\label{sec:ki-datenverarbeitung}

Der Einsatz von Künstlicher Intelligenz und maschinellem Lernen hat die Datenverarbeitung in den vergangenen Jahren grundlegend verändert. Schon vor dem Aufkommen großer Sprachmodelle wurden Verfahren des maschinellen Lernens eingesetzt, um Schritte der Datenaufbereitung, der Mustererkennung und der Entscheidungsunterstützung zu automatisieren und so aus großen Datenbeständen Erkenntnisse zu gewinnen \autocite{bharadiyaRoleMachineLearning2023,safitraAdvancementsArtificialIntelligence2024}. Insbesondere im Bereich der Datenaufbereitung haben sich lernende Verfahren als nützlich erwiesen, etwa bei der Datenbereinigung, der Imputation fehlender Werte, dem Abgleich von Entitäten und der Erkennung von Ausreißern \autocite{coteDataCleaningMachine2024,mumuniAutomatedDataProcessing2025}. Ein einschlägiges Beispiel ist HoloClean von Rekatsinas et al., das Integritätsbedingungen, externe Referenzdaten und statistische Eigenschaften des Bestands in einem probabilistischen Modell zusammenführt und daraus Reparaturvorschläge ableitet, statt für jeden Fehlertyp eine eigene Regel zu verlangen \autocite{rekatsinasHoloClean2017}. Bereits vor dem Aufkommen großer Sprachmodelle bestand der Beitrag lernender Verfahren damit vor allem darin, die vollständige Vorabformulierung der Bereinigungsvorschrift zu ersetzen, also genau jene Eigenschaft, die in dieser Arbeit für \acp{LLM} untersucht wird. Eine systematische Übersicht zeigt dabei ein wechselseitiges Verhältnis: Maschinelles Lernen wird einerseits zur Datenbereinigung genutzt, andererseits ist eine hohe Datenqualität Voraussetzung für die Leistungsfähigkeit der Modelle selbst \autocite{coteDataCleaningMachine2024,mumuniAutomatedDataProcessing2025}.

Mit der Verfügbarkeit dialogfähiger Sprachmodelle hat sich dieser Trend noch einmal verstärkt. Frühe Arbeiten beschreiben, wie Modelle wie ChatGPT Datenfachleute bei zahlreichen Teilschritten ihres Arbeitsablaufs unterstützen können, von der Datenbereinigung und Vorverarbeitung über die Modellbildung bis hin zur Interpretation der Ergebnisse \autocite{hassaniRoleChatGPTData2023,bharadiyaRoleMachineLearning2023}. Gleichzeitig betonen diese Beiträge, dass der Einsatz solcher Modelle mit Herausforderungen wie Verzerrungen, mangelnder Nachvollziehbarkeit und Fragen der Verlässlichkeit verbunden ist und daher einer kritischen Begleitung bedarf \autocite{hassaniRoleChatGPTData2023,safitraAdvancementsArtificialIntelligence2024}. Insgesamt zeichnet die Literatur das Bild einer zunehmend \ac{KI}-gestützten Datenverarbeitung, in der die Automatisierung von Routineaufgaben im Vordergrund steht, eine systematische empirische Bewertung der erzielten Qualität jedoch häufig noch aussteht \autocite{mumuniAutomatedDataProcessing2025,coteDataCleaningMachine2024}.

\section{LLMs für strukturierte Daten und synthetische Datensätze}
\label{sec:llm-tabellen}

Dieser Abschnitt klärt zwei Voraussetzungen der Untersuchung: dass \acp{LLM} tabellarische Daten überhaupt verarbeiten können, und dass ein synthetisch erzeugter Datensatz als Prüfgrundlage taugt.

Für den Einsatz von \acp{LLM} auf strukturierten Daten liegen direkte empirische Untersuchungen vor. Ein früher und vielzitierter Beitrag ist Ditto von Li et al.: Die Autoren führen den Abgleich zweier Datensätze auf eine Klassifikationsaufgabe zurück und lösen sie mit einem vortrainierten Sprachmodell, das für diesen Zweck nachtrainiert wird \autocite{liDeepEntityMatching2020}. Damit ist die in dieser Arbeit als anspruchsvollste Bereinigungsstufe geführte Erkennung unscharfer Duplikate bereits als originäres Anwendungsfeld vortrainierter Sprachmodelle belegt, allerdings mit aufgabenspezifisch angepassten Modellen und nicht, wie hier, mit unverändert eingesetzten. Narayan et al.\ gehen einen Schritt weiter und formulieren fünf Aufgaben der Datenbereinigung und -integration als reine Prompting-Aufgaben ohne Nachtraining, wobei sie neben der Ergebnisqualität auch die Empfindlichkeit gegenüber Prompt-Format und Beispielauswahl untersuchen \autocite{narayanFoundationModelsWrangle2022}. Damit ist die grundsätzliche Eignung bereits experimentell untersucht, die Ergebnisse sprechen zugleich dafür, konkrete Aufgaben und Prompt-Varianten getrennt zu bewerten.

Davon zu unterscheiden ist die Erzeugung synthetischer Daten durch generative Modelle. Übersichtsarbeiten behandeln deren Verfahren und Grenzen \autocite{longLLMsDrivenSyntheticData2024,fonsecaTabularLatentSpace2023,goyalSystematicReviewSynthetic2024}. Li et al.\ untersuchen beispielsweise LLM-generierte Trainingsdaten für Textklassifikation \autocite{liSyntheticDataGeneration2023}. Die vorliegende Arbeit erzeugt ihre Testdaten dagegen regelbasiert mit Faker. Ergebnisse über generierte Trainingsdaten sind daher kein unmittelbarer Qualitätsnachweis für diesen Testdatensatz. Dessen Eignung wird anhand des dokumentierten Datenmodells, der eingebrachten Fehler und der Validierung in Kapitel~\ref{ch:datensatz} begründet.

\section{Verwandte Arbeiten zur Automatisierung von ETL}
\label{sec:verwandte-arbeiten}

Die Automatisierung von \ac{ETL}-Prozessen mit Methoden der Künstlichen Intelligenz ist Gegenstand einer wachsenden Zahl von Arbeiten. Ein erster, schon länger verfolgter Ansatz nutzt klassische Verfahren des maschinellen Lernens, um einzelne Schritte der Pipeline zu automatisieren, etwa das Abbilden von Schemata, die Erkennung von Anomalien und die Ableitung von Transformationsregeln aus den Daten \autocite{mishraAutomatingDataIntegration2020,reddyAIDrivenDataIntegration2025}. Darauf aufbauend beschreibt eine Reihe neuerer Beiträge umfassendere Architekturen, in denen \ac{KI}-Komponenten auf allen Stufen des \ac{ETL}-Prozesses eingebettet werden, um Aufgaben wie automatisierte Schemaerkennung, intelligente Qualitätssicherung und Leistungsoptimierung zu übernehmen \autocite{walhaDataIntegrationTraditional2024,kasireddyRoleAIModern2026}.

Für die folgenden Beiträge gilt dabei der in Abschnitt~\ref{sec:recherche} benannte Vorbehalt in besonderem Maße: Sie beschreiben überwiegend Architekturentwürfe und Erwartungen, ohne diese empirisch zu prüfen, und erscheinen teilweise in Organen ohne strenge Begutachtung. Sie werden hier referiert, um den Erwartungshorizont abzustecken, gegen den die Ergebnisse dieser Arbeit in Abschnitt~\ref{subsec:einordnung} gestellt werden.

Ein wiederkehrendes Leitbild dieser Arbeiten ist die Vorstellung selbstheilender und adaptiver Datenpipelines, die Probleme proaktiv erkennen, eigenständig beheben und sich an veränderte Datenstrukturen anpassen, statt rein reaktiv auf Fehler zu reagieren \autocite{chakrabortyETLHowAI2025,reddyAIDrivenDataIntegration2025}. In diesem Zusammenhang werden zunehmend autonome \ac{KI}-Agenten diskutiert, die Datenaufnahme, Bereinigung und Integration weitgehend selbstständig ausführen sollen \autocite{maddaliAutonomousAIAgents2023,chakrabortyETLHowAI2025}. Mehrere Beiträge widmen sich dabei spezifischen Teilaufgaben, etwa der intelligenten Anomalieerkennung in \ac{ETL}-Strecken, der automatisierten Qualitätssicherung oder der automatisierten Erzeugung von Merkmalen \autocite{khanIntelligentAnomalyDetection2025,maddaliAutomatingDataQuality}. Insbesondere wird betont, dass nicht ein vollständiger Ersatz, sondern eine hybride Verknüpfung von regelbasierten Systemen und \ac{KI}-Modellen erstrebenswert sei, um Nachvollziehbarkeit und Kontrolle bei gleichzeitig höherer Flexibilität zu wahren \autocite{walhaDataIntegrationTraditional2024,kasireddyRoleAIModern2026}.

Ein weiterer empirischer Ansatz ist Jellyfish von Zhang et al.: Die Autoren passen lokale Sprachmodelle mit 7 bis 13 Milliarden Parametern durch Instruction-Tuning an Fehlererkennung, Imputation, Schemaabgleich und Entitätsabgleich an \autocite{zhangJellyfish2024}. Das verdeutlicht, dass neben Modellgröße und Betriebsform auch die aufgabenspezifische Anpassung zu berücksichtigen ist. Die hier untersuchte lokale Llama-Konfiguration bildet diese spezialisierte Modellklasse nicht ab.

Die vorliegende Untersuchung knüpft an diese Arbeiten an, evaluiert jedoch unverändert eingesetzte Modelle auf den selbst definierten Aufgaben. Ihr Schwerpunkt liegt auf der Gegenüberstellung mit einer regelbasierten Referenz und der gemeinsamen Betrachtung von Prompt-Vorlagen, Code-Generierung, Direktverarbeitung und verketteter Ausführung. Ein numerischer Leistungsvergleich mit den genannten Veröffentlichungen wäre wegen unterschiedlicher Datensätze, Aufgaben und Metriken nicht sachgerecht.

\section{Forschungslücke und Abgrenzung}
\label{sec:forschungsluecke}

Die Forschungslücke dieser Arbeit liegt nicht im grundsätzlichen Fehlen empirischer LLM-Evaluierungen. Die genannten Arbeiten untersuchen bereits mehrere Datenaufgaben und unterschiedliche Formen der Modellanpassung und Prompt-Gestaltung. Ergänzend wird hier für ein einheitlich konstruiertes \ac{ETL}-Szenario betrachtet, wie sich die Leistung zwischen unabhängigen Aufgaben, verketteter Verarbeitung und Durchläufen mit Soll-Zwischentabellen unterscheidet. Hinzu kommen der Vergleich von Code-Generierung und direkter Datenverarbeitung sowie die explorativen Ebenen zum Aufgabenzuschnitt.

Die wiederholte Ausführung dient dabei als querschnittliche Untersuchungsdimension. Sie beschreibt die beobachtete Schwankung der Kennzahlen unter den jeweiligen Versuchsbedingungen. Wegen der begrenzten Wiederholungszahl handelt es sich um eine deskriptive Untersuchung, eine allgemeine Aussage über die Reproduzierbarkeit einer Modellfamilie oder eine vollständige statistische Absicherung der Unterschiede wird damit nicht beansprucht.

FF1 adressiert die aufgabenbezogene Eignung gegenüber der regelbasierten Vergleichspipeline und die ergänzenden Ausführungsebenen. FF2 betrifft Unterschiede zwischen den drei konkret eingesetzten Prompt-Vorlagen. Die wiederholten Ausführungen helfen bei beiden Fragen, beobachtete Unterschiede im Verhältnis zur Streuung einzuordnen.